\section{Related work}\label{sec:related}

\paragraph{Equivariant interatomic potentials.}
SchNet \cite{schutt2017schnet} encodes atomic environments through invariant radial features, while DimeNet++ \cite{gasteiger2020dimenetpp} also incorporates angular information. NequIP \cite{batzner2022nequip} propagates equivariant tensor features, MACE \cite{batatia2022mace} builds higher-order correlations, and EquiformerV2 \cite{liao2024equiformerv2} combines equivariant attention with efficient eSCN convolutions \cite{passaro2023escn}. PaiNN \cite{schutt2021painn}, TorchMD-NET \cite{tholke2022torchmdnet} and ViSNet \cite{wang2024visnet} use scalar and vector features; FAENet \cite{duval2023faenet} uses stochastic frame averaging. We use these architectures as points of comparison for extending the Clifford substrate.

\paragraph{Cartesian tensor constructions.}
The equivalence between symmetric traceless Cartesian tensors and spherical irreducible representations is classical \cite{stone2013theory}; recent work provides orthonormal change-of-basis constructions up to rank 9 \cite{shao2025orthonormal}. ICTP \cite{zaverkin2024ictp} uses irreducible Cartesian tensor products, TensorNet \cite{simeon2023tensornet} uses rank-2 tensor operations, CACE \cite{cheng2024cace} formulates the atomic cluster expansion in Cartesian coordinates, and HotPP \cite{wang2024n} propagates tensors of arbitrary rank. GotenNet \cite{aykent2025gotennet} uses geometry-aware tensor attention, while Gaunt tensor products \cite{luo2024gaunt} offer efficient spherical-harmonic couplings. CliffordSTF draws on the same STF foundation as ICTP, but couples it to a parallel Clifford track through explicit generation and contraction maps (ESI Section~\ref{E-supp:ictp}).

\paragraph{Geometric algebra networks.}
Geometric Clifford algebra networks \cite{ruhe2023gcan}, Clifford group equivariant neural networks \cite{ruhe2023cgenn} and Clifford neural layers \cite{brandstetter2022clifford} use multivector representations and geometric products in neural networks. GATr \cite{brehmer2023geometric} uses projective geometric algebra, Clifford-steerable CNNs \cite{zhdanov2024clifford} construct kernels on pseudo-Riemannian manifolds, and Clifford group-equivariant simplicial message passing networks \cite{liu2024csmpn} extend Clifford message passing to simplicial complexes. Our focus is the explicit angular content of the $\mathrm{Cl}(3,0)$ substrate and its extension for interatomic potentials. OC20 reports force cosine similarity alongside force MAE \cite{chanussot2021oc20}. Prediction-error metrics alone can miss differences in molecular simulation behaviour \cite{kovacs2023mace}. We therefore assess directional agreement separately from force MAE on rMD17.
